\section{Discussion}
The document-paired validation strengthens the central TF-IDF finding: S+L+E removes 26.2\% of its gap, and entity masking improves cross-dataset AUC beyond matched random deletion both with and without length matching. The raw gap reduction is an operational property of these interventions on these corpora, not an estimate of the maximum or unique causal contribution of the cue families. Length matching changes the sampled population and uses class labels in the target corpus, so it remains a diagnostic rather than an unlabeled deployment method.

\emph{Random controls change the transformer interpretation.} Entity masking is worse than random deletion for DistilBERT in the paired comparison, rather than merely showing an inconclusive advantage. For RoBERTa, S+L+E removes about half the raw gap but has no clear pooled benefit over its length-matched random control. This distinction prevents a large raw gain from being attributed to the selected cues without accounting for nonspecific perturbation. The focused transformer design leaves source-only contrasts untested in the new run; their archived values should not be presented as newly validated findings.

\emph{Dataset pairs remain essential.} DistilBERT's negative residual-involving contrasts persist, but the analogous RoBERTa pattern is weaker and one interval is very wide. A pooled average conceals this uncertainty. The stronger claim that both transformers benefit only on ISOT--McIntire transfers is not supported consistently by the paired validation. Pair-level results remain exploratory rather than population-level evidence about dataset types.

\emph{External transfer is source dependent.} Celebrity provides a distinct target while retaining the source-side design. TF-IDF's primary entity-minus-random mean is positive for all sources, but the transformer response changes with the source corpus. In particular, the residual-trained RoBERTa advantage contrasts with its negative ISOT- and McIntire-trained entity means. The negative pooled DistilBERT entity result from the three-corpus analysis is not universal on the external target. These are meaningful boundary conditions, not grounds to discard adverse outcomes or pool them away. Because the target is small, collection-specific and only overlap-screened, the extension broadens the empirical setting without establishing generalization to all news domains.

\emph{Deletion changes transformer inputs.} A regularization-like response is one possible explanation for random deletion gains, but deletion also shifts later content into the first 256 tokens. We did not hold the original token window fixed, match subword counts, or isolate training-side from test-side changes. The results therefore support model-dependent sensitivity to these transformations, without identifying its mechanism.

\emph{The audit and the intervention answer different questions.} Normalized overlap checks motivated construction of the WELFake residual. Reuters markers support the interpretation of label direction but do not independently verify article veracity. The paired-validation audit checks saved sample identities and result provenance; it does not substitute for reproducing normalization, duplicate matching or human label validation from the source data.

For dataset-transfer studies, these findings support reporting overlap checks, label mappings, article-paired masking comparisons, matched random controls, model-specific responses and ordered dataset pairs. Length-matching comparisons should state explicitly that they alter the sampled population. Uncertain contrasts should be described as inconclusive rather than equivalent.
